%==============================================================================
\section{Act VI: The Proofreader (\texttt{LoopClosing})}
%==============================================================================

As the camera moves, small errors accumulate (drift). \texttt{LoopClosing} detects when the camera returns to a previously visited place and corrects the drift.

\subsection{Why Loop Closure Matters}

Imagine walking around a building and returning to your starting point:
\begin{itemize}
    \item Without loop closure: The map shows you're 10 meters away from the start
    \item With loop closure: The system recognizes you're back at start and fixes the error
\end{itemize}

\subsection{The Loop Closing Pipeline}

\begin{lstlisting}[caption={LoopClosing::Run() simplified}]
void LoopClosing::Run()
{
    while(!mbFinishRequested)
    {
        if(CheckNewKeyFrames())
        {
            // 1. Detect loop candidates using Bag-of-Words
            if(DetectLoop())
            {
                // 2. Compute similarity transformation (with scale)
                if(ComputeSim3())
                {
                    // 3. Correct the loop
                    CorrectLoop();
                    
                    // 4. Global Bundle Adjustment (in a new thread)
                    RunGlobalBundleAdjustment();
                }
            }
        }
    }
}
\end{lstlisting}

\subsection{Step 1: Loop Detection}

We use Bag-of-Words to find similar KeyFrames:

\begin{lstlisting}[caption={DetectLoop concept}]
bool LoopClosing::DetectLoop()
{
    // 1. Compute minimum score threshold
    //    (use similarity with neighbors as baseline)
    float minScore = ComputeMinScoreInNeighbors();
    
    // 2. Query database for similar KeyFrames
    vector<KeyFrame*> vpCandidates = 
        mpKeyFrameDB->DetectLoopCandidates(mpCurrentKF, minScore);
    
    if(vpCandidates.empty())
        return false;
    
    // 3. Check temporal consistency
    //    (candidates should appear in multiple consecutive KFs)
    for(KeyFrame* pKF : vpCandidates)
    {
        if(CheckConsistency(pKF))
            mvpEnoughConsistentCandidates.push_back(pKF);
    }
    
    return !mvpEnoughConsistentCandidates.empty();
}
\end{lstlisting}

\subsection{Step 2: Compute Sim3}

We need to find the transformation between current and loop KeyFrame:

\begin{lstlisting}[caption={ComputeSim3 concept}]
bool LoopClosing::ComputeSim3()
{
    for(KeyFrame* pKF : mvpEnoughConsistentCandidates)
    {
        // 1. Match features between current KF and candidate
        vector<MapPoint*> vpMatchedPoints;
        int nMatches = ORBmatcher::SearchByBoW(mpCurrentKF, pKF, 
                                               vpMatchedPoints);
        
        if(nMatches < 20)
            continue;
        
        // 2. Compute Sim3 using RANSAC
        Sim3Solver solver(mpCurrentKF, pKF, vpMatchedPoints);
        
        cv::Mat Scm;  // Similarity transformation
        if(solver.ComputeSim3(Scm))
        {
            mScw = Scm;
            mpMatchedKF = pKF;
            return true;
        }
    }
    return false;
}
\end{lstlisting}

\textbf{Note:} \texttt{Sim3} includes scale, which is important for monocular SLAM where scale can drift.

\subsection{Step 3: Correct Loop}

Once we find a valid loop, we correct the drift:

\begin{lstlisting}[caption={CorrectLoop concept}]
void LoopClosing::CorrectLoop()
{
    // 1. Stop LocalMapping temporarily
    mpLocalMapper->RequestStop();
    
    // 2. Compute correction for current KeyFrame
    Sophus::Sim3f Scw = mpCurrentKF->GetPose() * mScw.inverse();
    
    // 3. Propagate correction to all connected KeyFrames
    for(KeyFrame* pKF : mvpCurrentConnectedKFs)
    {
        Sophus::Sim3f CorrectedPose = Scw * pKF->GetPoseRelativeTo(mpCurrentKF);
        pKF->SetPose(CorrectedPose);
    }
    
    // 4. Correct all associated MapPoints
    for(MapPoint* pMP : mvpLoopMapPoints)
    {
        pMP->SetWorldPos(CorrectedPos);
    }
    
    // 5. Fuse duplicate MapPoints
    SearchAndFuse();
    
    // 6. Pose-graph optimization
    Optimizer::OptimizeEssentialGraph(...);
    
    // 7. Resume LocalMapping
    mpLocalMapper->Release();
}
\end{lstlisting}

\subsection{Step 4: Global Bundle Adjustment}

After loop closure, run global BA to refine everything:

\begin{lstlisting}[caption={Global BA}]
void LoopClosing::RunGlobalBundleAdjustment()
{
    // Run in a separate thread to not block loop closing
    thread* pThread = new thread(&LoopClosing::GlobalBundleAdjustmentThread, 
                                 this);
}

void LoopClosing::GlobalBundleAdjustmentThread()
{
    // Optimize ALL KeyFrames and MapPoints
    Optimizer::GlobalBundleAdjustment(mpAtlas->GetCurrentMap());
    
    // Merge corrected map back
    // ...
}
\end{lstlisting}

\subsection{The Covisibility Graph}

Loop closing uses the covisibility graph to propagate corrections:

\begin{itemize}
    \item \textbf{Nodes}: KeyFrames
    \item \textbf{Edges}: Weighted by number of shared MapPoints
    \item \textbf{Use}: Find which KeyFrames are affected by a correction
\end{itemize}

\begin{lstlisting}[caption={Covisibility example}]
// Get KeyFrames that share at least 15 MapPoints with pKF
vector<KeyFrame*> vpCovisibleKFs = pKF->GetCovisiblesByWeight(15);
\end{lstlisting}
